\chapter{Derivatives of functions $\R \to \R$}

Throughout this chapter, let $U$ denote an open subset of $\R$. We'll formalize the theory of derivatives of functions $U \to \R$. The key idea is that a derivative represents the best possible linear approximation of a function.

\section{Slope of a tangent line}

One way of formalizing the idea of a best linear approximation 

\begin{definition}
	A function $f : U \to \R$ is \emph{differentiable} at a point $a \in U$ if 
	\[ \lim_{h \to 0} \frac{f(a+h) - f(a)}{h} \]
	exists. If this limit does exist, its value is called the \emph{derivative of $f$ at $a$}, and is denoted $f'(a)$. 
\end{definition}

\section{Best linear approximation}

\begin{proposition}
	Suppose $f : U \to \R$ is differentiable at $a \in U$ if and only if there exists a linear map $\ell : \R \to \R$ such that
	\[ f(a+h)-f(a) - \ell(h) = o(|h|) \text{ as } h \to 0. \]
	Moreover, if there exists such an $\ell$, then we must have $\ell(h) = f'(a)h$ for all $h \in \R$. 
\end{proposition}

\begin{definition}
	If $f : U \to \R$ is differentiable at a point $a \in U$, we define $df_a : \R \to \R$ to be the linear map $df_a(h) = f'(a)h$. 
\end{definition}